% Discussion for item 1.
% scripts/make_latex.py will NEVER overwrite this file -- edit it freely.

The two outputs differ in the order of the last two lines: in 1.1 the parent
printed first, in 1.2 the child did.

In 1.1 the parent does not wait. Once \texttt{fork()} returns there are two
runnable processes, and nothing in the program orders them --- which one reaches
its \texttt{printf} first is decided by the operating system's scheduler. A
single run cannot show that, since it only ever shows one of the possible
orders, so we ran 1.1 repeatedly and counted. The child does come first
sometimes, but rarely: once in 500 runs on an idle machine. Loading the machine
with eight busy loops raised it to 44 in 500. The order really is undecided ---
it is simply that on an idle multi-core machine the parent almost always wins
the race, and making the processes compete for the CPU is what lets the child
win often enough to see.

In 1.2 the parent calls \texttt{wait(NULL)}, which blocks it until the child
terminates. The parent therefore cannot reach its \texttt{printf} until the
child has already run and exited, so the child's line always comes first ---
under load or not. Adding \texttt{wait()} turns the ordering from an accident of
scheduling into a property the program guarantees.
